\appendixsection{A CONNECTEDNESS PROPERTY}

Lemma 1. Let X be a connected topological space and \(\Omega\) a dense open subset of X. If, for any \(x \in X\) , there exists a neighbourhood V of x such that \(V \cap \Omega\) is connected, then \(\Omega\) is connected.

Indeed, let \(\Omega_0\) be a non-empty open and closed subset of \(\Omega\) . Let \(x \in X\) and let \(V\) be a neighbourhood of \(x\) such that \(V \cap \Omega\) is connected. If \(x \in \overline{\Omega}_0\) ,

\[
(\mathrm{V} \cap \Omega) \cap \Omega_ {0} = \mathrm{V} \cap \Omega_ {0} \neq \varnothing ,
\]

so \(V \cap \Omega \subset \Omega_{0}\) . Thus, since \(\Omega\) is dense in X, \(\overline{\Omega}_{0}\) is a neighbourhood of x. Consequently, \(\overline{\Omega}_{0}\) is non-empty, open and closed, and since X is connected, \(\overline{\Omega}_{0} = X\) . Since \(\Omega_{0}\) is closed in \(\Omega\) , this implies that \(\Omega_{0} = \Omega \cap \overline{\Omega}_{0} = \Omega\) , which proves that \(\Omega\) is connected.

Lemma 2. Let U be an open ball in \(\mathbf{C}^n\) and \(f: \mathrm{U} \to \mathbf{C}\) a holomorphic function, not identically zero. Let A be a subset of U such that \(f = 0\) on A. Then U-A is dense in U and connected.

The density of U-A follows from Differentiable and Analytic Manifolds, Results, 3.2.5. Assume first that \(n = 1\) . If \(a \in \mathrm{A}\) , the power series expansion of \(f\) about \(a\) (Differentiable and Analytic Manifolds, Results, 3.2.1) is not reduced to 0, and it follows that there exists a neighbourhood \(\mathrm{V}_a\) of \(a\) in U such that \(f\) does not vanish on \(\mathrm{V}_a - \{a\}\) . Thus, \(a\) is isolated in A, which proves that A is a discrete subset of U, hence countable since U is countable at infinity. Let \(x, y \in \mathrm{U} - \mathrm{A}\) . The union of the real affine lines joining \(x\) (resp. \(y\) ) to a point of A is meagre (General Topology, Chap. IX, §5, no. 2, Def. 2). Hence, there exists \(z \in \mathrm{U} - \mathrm{A}\) such that neither of the segments \([x, z]\) and \([y, z]\) meets A. The points \(x, y, z\) thus belong to the same connected component of U-A, which proves the lemma in the case \(n = 1\) . We turn to the general case. We can assume that A is the set of zeros of \(f\) (General Topology, Chap. I, §11, no. 1, Prop. 1). Let \(x, y \in \mathrm{U} - \mathrm{A}\) and let L be an affine line containing \(x\) and \(y\) . The restriction of \(f\) to \(\mathrm{L} \cap \mathrm{U}\) is not identically zero since \(x \in \mathrm{L} \cap \mathrm{U}\) . By what has already been proved, \(x\) and \(y\) belong to the same connected component of \((\mathrm{L} \cap \mathrm{U}) - (\mathrm{L} \cap \mathrm{A})\) and hence to the same connected component of U-A.

Lemma 3. Let X be a finite dimensional connected complex-analytic manifold and let A be a subset of X satisfying the following condition:

For any \(x \in X\) , there exists an analytic function germ \(f_{x}\) , not vanishing at x, such that the germ of A at x is contained in the germ at x of the set of zeros of \(f_{x}\) .

Then X-A is dense in X and connected.

The density of X-A follows from Differentiable and Analytic Manifolds, Results, 3.2.5. We can assume that A is closed (General Topology, Chap. I, §11, no. 1, Prop. 1). For any \(x \in \mathrm{X}\) , there exists an open neighbourhood \(\mathrm{V}\) of \(x\) and an isomorphism \(c\) from \(\mathrm{V}\) to an open ball in \(\mathbf{C}^n\) such that \(c(\mathrm{A} \cap \mathrm{V})\) is contained in the set of zeros of a holomorphic function not identically zero on \(c(\mathrm{V})\) . Then, by Lemma 2, \(\mathrm{V} \cap (\mathrm{X} - \mathrm{A})\) is connected. In view of Lemma 1, this proves that \(\mathrm{X} - \mathrm{A}\) is connected.

